How to run PythonTeX code?

**Compiling a document that uses PythonTEX involves three steps: run LATEX, run pythontex.py, and finally run LATEX again.
You may wish to create a symlink or launching wrapper for pythontex.py, if one was not created during installation.**

 - Generally, PythonTeX offers 7 inline commands to be used within nor-
mal LaTeX running text:
– \py Evaluates a Python expression and prints its value. Here
also functions defined via pycode can be used.
– \pyc Executes code, output goes to stdout, does not typeset it
– \pyb Executes code, but outputs only the code prettyprinted.
To show anything printed, use \printpythontex or \stdoutpythontex.
– \pys Supports variable and expression substitution
– \pyv Prettyprints code (nothing is executed).
- Same as \pygment{python}
- Other than the pyverbatim env no new paragraph is created.
– \pygment General code typesetting
– \pycon Accesses a variable from a previous pyconsole block
